% ======================================================================
\section{Conclusion}\label{sec:conclusion}

A deaf viewer misses the sounds whose source is not on screen, and no existing channel supplies them automatically.
This project built a system that does, from open models with no training, together with a labelled benchmark and a
metric that prices what a viewer loses. On the test set the final system finds 37\,\% of the needed sounds, half of its
pictures are right, and its cost is significantly lower than showing nothing ($p = 0.047$) and lower than the September
version. The remaining errors sit where today's open models stop: sounds the detectors barely hear under speech and
music, and a vision model that takes ``the object is there'' for ``the object is making the sound''.

\paragraph{Answers to the research questions.}
\begin{description}[style=nextline]
  \item[RQ1 -- which audio information matters most?] Discrete, sourced events whose source is off screen: sirens,
        alarms, bells, a baby crying, glass, a knock. The benchmark's importance rule (an event one can say in a
        sentence, or danger) captures them; steady place noise (wind, traffic hum) is excluded by design, and it is
        also where detectors, gate and labels disagree most. Whether DHH viewers rank sounds the same way is untested.
  \item[RQ2 -- what granularity?] Decide and score at the level of the AudioSet \emph{family} (depth $\geq 1$), draw the
        specific source. Family matching was the one de-duplication and scoring rule that generalised without a
        threshold; the picture names the maker (a barking dog, not ``bark'') from a hand table, which beat generic VLM
        wording 7/7 to 3/7.
  \item[RQ3 -- can generated augmentations improve accessibility beyond subtitles?] Not shown. What is shown is that
        the system adds information no subtitle carries (off-screen sounds) at a modelled cost below showing nothing,
        and that its pictures are recognised at a glance about half the time. A study with viewers is the necessary next
        step.
  \item[RQ4 -- how to evaluate?] Per sound, against human labels, with the error split by cause and a cost that prices
        a miss against a wrong picture, reported over a range of weights; decisions pre-registered on a development set.
        Automatic protocols with model-written references or model judges failed: they are circular, decide visibility
        at chance, or cannot tell a picture from a text tag.
  \item[RQ5 -- comparison with audio-to-visual approaches?] Audio-to-image models draw the whole scene and cannot leave
        a visible object out, so they were represented by a blind baseline with the same detector and generator. Against
        it, the gate removes about half a wrong picture per clip at a small cost in recall (Holm-significant on the
        September test table); F1 does not move because the gate's value is in what it does not show.
\end{description}
